\chapter{ОПИСАНИЕ ПРЕДЛАГАЕМОГО РЕШЕНИЯ}

\section{Архитектура продакшн-решения и System Design}

В первых двух главах были рассмотрены предметная область, архитектурные ограничения исходной системы и дизайн экспериментальной проверки предлагаемых доработок. В третьей главе внимание переносится с уровня проектирования на уровень реализации: описывается целевая архитектура продакшн-решения, принципы интеграции ClawdLab и OpenClaw, особенности развертывания на собственной инфраструктуре и подходы к сопровождению системы.

С практической точки зрения разрабатываемое решение представляет собой self-hosted агентную среду для исследовательских задач, развернутую на VPS-сервере под управлением Ubuntu. В качестве технологической основы используются OpenClaw как инфраструктурный слой подключения и исполнения агентов и ClawdLab как лабораторная надстройка, организующая исследовательский процесс, роли агентов и выполнение задач. Такое разделение позволяет отдельно рассматривать коммуникационный и интеграционный слой, прикладную исследовательскую логику, а также механизмы памяти, цикличности и воспроизводимости.

\subsection{Цели и требования к целевой архитектуре}

В контексте данной работы целевая архитектура должна удовлетворять нескольким требованиям. Она должна быть воспроизводимой и управляемой при развертывании на выделенном сервере без зависимости от внешней облачной платформы оркестрации. И она должна поддерживать независимый запуск нескольких агентных контуров, поскольку OpenClaw и ClawdLab выполняют разные функции и могут обслуживаться раздельно. Также архитектура должна допускать локальные доработки исходного кода, необходимые для адаптации системы под исследовательский сценарий.

Дополнительно к этим требованиям предъявляются требования к прозрачности выполнения задач, расширяемости провайдеров данных, контролю над конфигурацией окружения и удобству сопровождения. Поскольку работа ориентирована не только на запуск MVP, но и на дальнейшее развитие решения, архитектура должна допускать поэтапное внедрение доработок, в том числе механизмов структурированной памяти и управляемой цикличности исполнения.

\subsection{Логическая структура решения}

На логическом уровне систему можно представить как набор взаимосвязанных слоев. Нижний слой образует серверная инфраструктура: VPS-сервер, операционная система Ubuntu, сетевые настройки, переменные окружения и управляющие сессии. Поверх него располагается слой исполнения приложений, в котором OpenClaw и ClawdLab запускаются как два отдельных процесса в разных tmux-сессиях. Такое решение упрощает изоляцию, перезапуск и наблюдение за журналами каждого компонента.

Следующий слой составляет прикладная логика ClawdLab: маршруты API, постановка и запуск задач, взаимодействие с провайдерами, обработка результатов и формирование исследовательских артефактов.

Над этим слоем располагаются механизмы исследовательского управления, включающие память, фиксацию артефактов, логику повторных попыток и, в перспективе, управляемую цикличность. 

\subsection{Физическое развертывание на VPS-сервере}

Физическая схема развертывания в текущей реализации является минималистичной и ориентированной на MVP. Все основные компоненты размещаются на одном VPS-сервере под управлением Ubuntu. OpenClaw и ClawdLab запускаются вручную в отдельных tmux-сессиях, что позволяет независимо управлять жизненным циклом каждого сервиса, просматривать их вывод и при необходимости быстро восстанавливать рабочее состояние после сбоев или обновлений.

Такой способ развертывания не является финальным, но он соответствует этапу прототипирования. Он обеспечивает достаточный уровень контроля над окружением, прозрачность исполнения и низкий порог внесения изменений в кодовую базу. В дальнейшем эта схема может быть расширена до использования systemd, контейнеризации или внешнего reverse proxy, но на текущем этапе ручное управление через tmux лучше согласуется с задачами отладки и быстрой итерации.

\section{Разработка сервиса и MVP-интеграция}

В рамках MVP основное внимание уделяется не полному перепроектированию ClawdLab, а обеспечению работоспособного исследовательского контура, который можно развернуть на собственном сервере, настроить через переменные окружения и использовать для выполнения задач обзора литературы и анализа.

\subsection{Базовый контур MVP}

\subsection{Доработка literature provider}

В ходе запуска сервиса было выявлено, что система не поддерживает подключение к провайдерам литературы напрямую, для этого нужна надстройка, которая позволила бы агентам правильно обращаться в источниики данных. Для этого в файл lib/providers.ts была добавлена логика работы с источниками, предлагаемыми авторами clawdlab (пока что только europe pmc). Эта доработка позволяет использовать Europe PMC fallback как внутренний источник литературы.

Вторым изменением стало обновление маршрута запуска literature job: логика маршрута была адаптирована так, чтобы задача могла завершаться сразу, без необходимости ожидания внешнего асинхронного провайдера. Тем самым был упрощен путь исполнения запроса в конфигурации, где используется встроенный провайдер Europe PMC.

С инженерной точки зрения это изменение уменьшает связность между стартом задачи и внешним асинхронным контуром выполнения. Для MVP это особенно полезно, поскольку снижает сложность интеграции, упрощает отладку и уменьшает число состояний, которые необходимо учитывать при сопровождении системы. Одновременно такое изменение формирует более чистую основу для дальнейшего введения нескольких режимов исполнения: синхронного, асинхронного и гибридного.

Чтобы сделанные доработки могли использоваться не только в локальной рабочей копии, но и в повторяемом процессе развертывания, были обновлены конфигурационные и пользовательские документы. Новые переменные окружения были задокументированы в .env.example и README.md. Это позволяет явно зафиксировать поддерживаемый режим запуска literature provider и сократить число скрытых зависимостей при переносе решения на другой сервер или при повторной установке.

Фактически это означает, что сценарий получения литературы может быть активирован без отдельного внешнего literature agent API.

\subsection{Архитектура модуля структурированной памяти}

В данной работе внедряется следующий вариант HCC-памяти, ориентированный на запуск слоя исследовательской памяти. В отличие от системы ML Masters 2.0 в этой версии не вводятся отдельный слой снимков контекста (context snapshot), графовая модель связей, поля sourceVersion и интервалы актуальности validFrom/validTo. Такое ограничение уменьшает сложность схемы данных и логики сопровождения, но при этом сохраняет главный архитектурный принцип HCC: разделение памяти по уровням зрелости и роли в исследовательском процессе.

В основе модуля лежат три сущности: MemoryRecord, MemorySourceRef и MemoryPromotionJob. Вместе они образуют минимально достаточную архитектуру для хранения самих записей памяти, фиксации происхождения знаний и наблюдаемого продвижения памяти между уровнями. Тем самым достигается полноценная трехуровневая модель, включающая рабочую, эпизодическую и каноническую память, но без преждевременного усложнения модели графами и версиями контекста.

MemoryRecord выступает главной сущностью памяти. Эта таблица хранит принадлежность записи лаборатории и, при необходимости, конкретной задаче или состоянию лаборатории, уровень памяти (working, episodic, canonical), тип записи, статус, текстовое содержимое, а также дополнительные структурированные данные. Для поддержки практического использования в записи также фиксируются уверенность, важность, источник происхождения, авторство и связи эволюции памяти. Последнее реализуется через три типа lineage-ссылок: parentMemoryId для иерархического родительства, promotedFromId для указания записи, из которой текущая память была поднята на следующий уровень, и supersedesId для замещения более старой версии знания новой.

Такое устройство MemoryRecord позволяет реализовать HCC-логику даже без графовой модели. Поле уровня задает архитектурное разделение памяти, тип записи нужен для извлечения релевантного контекста и фильтрации, а поля sourceType/sourceId и lineage-отношения обеспечивают происхождение данных (provenance) и прослеживаемость эволюции знания. Иначе говоря, система получает возможность не только хранить вывод, но и понимать, откуда он возник, к какому эпизоду относится и какую запись он уточняет или заменяет.

MemorySourceRef используется как дополнительный слой ссылок на источники. Он нужен потому, что одна запись памяти в реальном исследовательском процессе нередко опирается сразу на несколько сущностей. Например, итоговая запись типа canonical/finding может быть выведена не из одного объекта, а из комбинации задачи, голосования, документа, результата задания внешнего провайдера и лабораторного состояния. Поэтому вместо жесткой графовой модели вводится нормализованный слой происхождения данных (provenance), в котором для каждой записи памяти можно указать тип источника, его идентификатор, краткий ярлык, фрагмент источника, указатель на место в источнике и произвольные метаданные.

MemoryPromotionJob представляет собой журнал дистилляции и продвижения памяти между уровнями. Он хранит, в какой лаборатории и для какой задачи или состояния запускалось продвижение, из какого уровня в какой оно выполнялось, какой триггер его вызвал, какие записи были использованы на входе, какие были созданы на выходе, а также итоговый статус, краткую сводку и возможную ошибку. Наличие такой сущности принципиально важно, поскольку без нее HCC-продвижение превращается в скрытую внутреннюю логику. С журналом же процесс дистилляции становится наблюдаемым, воспроизводимым и пригодным для отладки.

На уровне логики памяти первый релиз поддерживает три слоя. Рабочая память играет роль краткосрочного рабочего контекста. В нее попадают текущий контекст задачи, активный план, свежие результаты провайдеров, временные заметки по исследовательской ветке и недавние артефакты, которые еще не прошли осмысление. Для этого уровня характерны быстрое обновление, допустимость шума и незавершенности, а также возможность истечения по expiresAt.

Эпизодическая память фиксирует завершенные или частично завершенные эпизоды работы. Сюда относятся краткие итоги попыток, локальные ошибки и обходные решения, сводки критики, сводки обсуждений, извлеченные уроки по отдельным эпизодам и промежуточные выводы, которые еще не должны считаться каноном. По сравнению с рабочим слоем этот слой уже сжат и структурирован, но все еще может содержать конкурирующие или конфликтующие версии интерпретации результатов.

Каноническая память хранит устойчивое знание лаборатории. На этот уровень поднимаются принятые выводы, сводки состояния лаборатории, финальные выводы, устойчивые правила лаборатории и межзадачные уроки, которые действительно подтверждены. В отличие от рабочей и эпизодической памяти, канонический слой изменяется редко, обычно не истекает по времени и обновляется не перезаписью, а через создание новой записи с указанием supersedesId. Это важно для сохранения истории эволюции знаний и предотвращения потери старых, но когда-то значимых выводов.

Рабочие записи создаются, когда задача взята в работу, запущено или завершено задание внешнего провайдера, появился новый активный план либо был загружен новый набор данных или артефакт. Эпизодические записи создаются, когда задача завершена, появляется критика, ветка обсуждения дошла до состояния, пригодного для сводки, или в одной ветке накопилось достаточно активности. Канонические записи создаются тогда, когда задача перешла в статус принятой, состояние лаборатории активировано или завершено, итоговый документ признан финальным либо PI или задание консолидации поднимает устойчивый вывод. 

Базовый пайплайн продвижения памяти между уровнями строится как трехшаговая схема. Первый шаг - переход из рабочей памяти в эпизодическую память. Он выполняется тогда, когда активная работа завершилась и из набора сырых результатов уже можно собрать короткий осмысленный эпизод. Второй шаг - переход из эпизодической памяти в каноническую память. Он выполняется, когда результат прошел через принятие, голосование или решение о состоянии, вывод повторяется и подтверждается или вошел в лабораторные документы и устойчивые выводы. Третий шаг - замещение внутри канонической памяти. Он нужен в тех случаях, когда старый вывод уточняется, новая версия надежнее и предыдущая запись больше не должна считаться актуальной. Тогда новая запись не затирает старую, а получает ссылку supersedesId на прежний finding.

Чтение памяти также строится по правилам извлечения релевантного контекста. Для активной задачи сначала извлекаются рабочие записи по taskId, затем эпизодические записи по той же задаче, а затем релевантные канонические записи по labId и нужным типам памяти. 

С точки зрения интеграции с текущей моделью ClawdLab новая память не заменяет существующие сущности, а располагается поверх них. Task и ProviderJob поставляют данные для рабочего слоя, TaskCritique, LabDiscussion и LabActivityLog образуют основу эпизодического уровня, а принятая задача, LabState и LabDocument становятся естественными источниками канонических знаний. Тем самым архитектура памяти остается согласованной с уже существующей доменной моделью и не требует разрушительного переписывания текущего приложения.

Следующим шагом после фиксации этой схемы является описание API, т. е. маршрутов вида POST /api/labs/slug/memory, GET /api/labs/slug/memory, и т. д. Параллельно с этим определяютсяавтоматические триггеры запуска продвижения: после завершения задачи, после перехода голосования в статус accepted/rejected, после финализации документа и после завершения состояния лаборатории. Именно такое поэтапное введение API и триггеров делает HCC-память полезной доработкой для данной системы.Ы


\subsection{Архитектура доработки цикличности исполнения}

 В текущей логике ClawdLab элементы цикла уже возможны: агент или PI может создать новую задачу, оставить критику, вернуться к обсуждению или попросить доработать результат. Однако такая цикличность остается неформальной: система не фиксирует, почему был начат новый проход, какая ошибка исправляется, чем новая попытка отличается от предыдущей и при каком условии цикл должен завершиться. Поэтому архитектурно предлагается отделить задачу как цель исследования от попытки выполнения как конкретного прохода решения. Для этого может быть введена сущность TaskAttempt или TaskRun, связанная с исходной задачей, результатом критики, использованными артефактами, статусом попытки и причиной повторного запуска.

Фиксированные правила цикличности могут быть описаны через отдельную политику выполнения для типа задачи или конкретной лаборатории. Такая политика задает допустимые триггеры повтора: наличие критики со статусом "нужно исправить", обнаружение пропущенных источников, неподтвержденные числовые значения, противоречие между полями результата или недостаточная трассируемость вывода к источникам. В этой же политике задаются условия остановки: максимальное число попыток, отсутствие улучшения по выбранным метрикам, превышение бюджета времени или стоимости, а также ручное решение PI принять текущий результат. Тогда цикл выполнения становится не скрытым поведением агента, а наблюдаемым конечным процессом: задача порождает попытку, попытка получает результат и критику, правило принимает решение о завершении или повторе, а новая попытка наследует только релевантный контекст и явно указывает, какую проблему она должна исправить.

С точки зрения реализации этот слой встроен в существующую серверную логику ClawdLab без изменения роли внешних агентов. Сервер хранит состояние цикла, связывает TaskAttempt с Task, TaskCritique, MemoryRecord и артефактами, а внешний агент получает уже подготовленный контекст следующей попытки. API включает операции создания попытки, завершения попытки, регистрации критики и принятия решения о следующем шаге. Такой подход сохраняет ClawdLab как координационный слой и одновременно делает цикличность проверяемой, воспроизводимой и пригодной для последующего сравнения в эксперименте.

\section{Развертывание, DevOps и MLOps-практики}

\subsection{Организация развертывания и запуск сервисов}

\section{Тестирование и верификация решения}

\subsection{План тестирования и валидации}
